\chapter{Devices \& Analysis}

A device is a physical component that can be used to build circuits.
Devices can be classified based on their functionality, characteristics,
and behavior. Here we will classify devices based on if they require
an external power source to operate or not, which leads us to two
categories of devices:
\begin{itemize}
    \item \textbf{Passive Devices}: These devices do not require a power
      source to operate. They can attenuate, store and dissipate
      electromagnetic energy.
    \item \textbf{Active Devices}: These devices require a power source to
      operate. They can attenuate, store, dissipate \textbf{and amplify}
      electromagnetic energy.
\end{itemize}

There are several important terms we will use throughout this handbook
and it is good to be familiar with them early on.
\begin{enumerate}
    \item \textbf{Port/Terminal}: A point of connection for a device.
      It can be a pin, contact, or lead. Two Terminal and Two Port
      devices are the most common.
    \item \textbf{Linear}: The relationship between two quantities (here
      voltage and current) is proportional by some factor. If you plot
      this relationship, you will get a straight line. If the relationship
      is not proportional, it is non-linear.
    \item \textbf{Time Invariant}: If a quantity does not change with time,
      it is time invariant. Many devices suffer from aging and therefore
      are time variant over long periods. For everyday circuits, we can
      assume that devices are time invariant.
    \item \textbf{Dissipation}: The act of releasing energy (usually in the
      form of heat) is known as dissipation.
\end{enumerate}

\section{Resistor}
An ideal resistor is a two-port, passive, linear and time-invariant
device that opposes the flow of current.

\begin{itemize}
  \item \textbf{Resistance:}
    This act of opposition is called resistance and is measured in
    Ohms (\(\Omega\)). A good way to imagine resistance is to think
    about its value as a proportionality factor between voltage and
    current. In fact, This is precisely the relationship that gives
    us Ohm's law:
    \[
      V = I_ \cdot R
    \]
  \item \textbf{Conductance:}
    The inverse of resistance is called conductance (\(G\)) and is
    measured in Siemens (S). The relationship between resistance and
    conductance is:
    \[
      G = \frac{1}{R}
    \]
    We'll mostly use conductance \(G\) when we analyse resistors
    mathematically as it simplifies the analysis a lot. The rearranged
    Ohm's law in terms of conductance is:
    \[
      I = V \cdot G
    \]
  \item \textbf{Power Dissipation:}
    A resistor dissipates energy in the form of heat. The power
    dissipated by a resistor is given by:
    \[
      P = I^2 \cdot R = V^2 \cdot G
    \]
  \item \textbf{Non-Ideal Behavior:}
    In practise, resistor's are not ideal and thus most of our
    assumptions above (linear, time-invariant, dissipative) fall
    apart at very high frequencies. However for the most part,
    we can treat resistors as ideal devices.
  \item \textbf{Circuit Analysis:}
    The behaviour of resistor is not affected by the kind of Input
    Signal (AC or DC) applied to it. Therefore, The same mathematical
    model can be used to describe the behaviour of a resistor in both
    AC and DC inputs.
\end{itemize}

\subsection{Capacitor}
A capacitor is a two-port passive device that stores energy in the form of
an electric field. The property of a capacitor is called capacitance and is
measured in Farads (F). The core relationship between charge, capacitance
and voltage is given by:
\begin{align*}
  Q = C \cdot V \\
  \frac{dQ}{dt} &= C \cdot \frac{dV}{dt} \hspace{1cm} \text{(assuming \(C\) is constant)} \\
  I &= C \cdot \frac{dV}{dt} \hspace{1cm} [I = \frac{dQ}{dt}]
\end{align*}
This is the fundamental relationship between current and voltage in a
capacitor or if you may call this the capacitor's version of Ohm's law.
The inverse of capacitance is called elastance (\(S\)) and is measured in Farads
inverse (F\(^{-1}\)).

The impedance of a capacitor is given by:
\[
  Z_C = \frac{1}{j\omega C} = \frac{1}{j 2 \pi f C} = \frac{1}{s C}
\]
The inverse of impedance is called admittance (\(Y\)) and is measured in
Siemens (S). The relationship between impedance and admittance is:
\[
  Y = \frac{1}{Z} = \frac{1}{\frac{1}{sC}} = sC
\]
You'll notice when we apply KCL to circuits within this book, we'll mostly
use admittance \(Y\) to express the impedance, we do this as it simplifies
the analysis a lot. The admittance version of Ohm's law is:
\[
  I = V \cdot Y
\]
The impedance of a capacitor is inversely proportional to frequency. This
means that at low frequencies, the impedance of a capacitor is very high and
at frequencies, the impedance of a capacitor is very low. This is why
capacitor effects are more pronounced at high frequencies as their
impedance drops enough to load the circuit. This is also why capacitors are
used for decoupling (they act as a short circuit at high frequencies).

An ideal capacitor is a linear (that is, the relationship between voltage and
current is a straight line) and time-invariant (that is, the relationship
between voltage and current does not change with time) of the three passive
components. It is also the only device that in its ideal form does not
dissipate energy.

A capacitor can be formed anywhere two conductors are separated by an
insulator. The charges are stored on the surface of the  conductors and the
energy is stored in the electric field between the  conductors. The energy
stored in a capacitor is given by:
\begin{align*}
  I &= C \cdot \frac{dV}{dt} \\
  P &= V \cdot I \\
  P &= V \cdot C \cdot \frac{dV}{dt} \\
  P &= C \cdot V \cdot \frac{dV}{dt} \\
  P &= C \cdot \frac{1}{2} \cdot \frac{d(V^2)}{dt} \hspace{1cm} [\frac{d(V^2)}{dt} = 2V \cdot \frac{dV}{dt}] \\
  \int P dt &= \int C \cdot \frac{1}{2} \cdot \frac{d(V^2)}{dt} dt \\
  E &= \frac{1}{2} C V^2
\end{align*}

\subsection{Inductor}
An ideal inductor is a two-port, passive, linear, time-invariant
and dissipative device that stores energy in the form of a
magnetic field.

\begin{itemize}
  \item \textbf{Inductance:}
    The relationship between flux and current is
    expressed by a proportionality factor called inductance (\(L\))
    and is measured in Henrys (H). It also gives us the inductor's
    version of Ohm's law:
    \begin{align*}
      \Phi &= L \cdot I \\
      \text{Differentiating both sides with respect to time:} \\
      \frac{d\Phi}{dt} &= L \cdot \frac{dI}{dt} \hspace{1cm} \text{(assuming \(L\) is constant)} \\
      V &= L \cdot \frac{dI}{dt} \hspace{1cm} [V = \frac{d\Phi}{dt}]
    \end{align*}
    The inverse of inductance is called reluctance (\(R\)) and is
    measured in Henrys inverse (H\(^{-1}\)).
    \[
      R = \frac{1}{L}
    \]
  \item \textbf{Impedance:}
    In the same way we did for capacitors, we can use Laplace
    Transform to  define the impedance of an inductor as the ratio
    of voltage to current in the Laplace domain:
    \begin{align*}
      V &= L \cdot \frac{dI}{dt} \\
      V(s) &= L \cdot s (I(s) - I(0)) \hspace{1cm} [\text{Laplace Transform}] \\
      V(s) &= s L I(s) \hspace{1cm} [\text{assuming \(I(0) = 0\)}]
    \end{align*}
    We can now define the impedance of an inductor as:
    \[
      Z_L = s L = j\omega L = j 2 \pi f L
    \]
  \item \textbf{Admittance:}
    The inverse of impedance is called admittance (\(Y\)) and is
    measured in Siemens (S). The relationship between impedance
    and admittance is:
    \[
      Y = \frac{1}{Z} = \frac{1}{sL}
    \]
    You'll notice when we apply KCL to circuits within this book,
    we'll mostly use admittance \(Y\) to express the impedance, we
    do this as it simplifies the analysis a lot. The admittance
    version of Ohm's law is:
    \[
      I = V \cdot Y
    \]
  \item \textbf{Frequency Response:}
    The impedance of an inductor is directly proportional to
    frequency. At low frequencies, the impedance of an inductor is
    very low and at high frequencies, the impedance of an inductor
    is very high. This means that at low frequencies, the inductor
    does not load the circuit which is often a very desirable
    effect  in many applications. This is also why inductors are
    used for filtering/chokes (they act as a short circuit at low
    frequencies but as an open circuit at high frequencies).
  \item \textbf{Energy and Power:}
    The energy and power stored in an inductor can be derived
    from the voltage-current relationship of an inductor:
    \begin{align*}
      V &= L \cdot \frac{dI}{dt} \\
      P &= V \cdot I \\
      P &= L \cdot \frac{dI}{dt} \cdot I \\
      P &= L \cdot I \cdot \frac{dI}{dt} \\
      P &= L \cdot \frac{1}{2} \cdot \frac{d(I^2)}{dt} \hspace{1cm} [\frac{d(I^2)}{dt} = 2I \cdot \frac{dI}{dt}] \\
      \int P dt &= \int L \cdot \frac{1}{2} \cdot \frac{d(I^2)}{dt} dt \\
      E &= \frac{1}{2} L I^2
    \end{align*}
  \item \textbf{Eddy Current Losses:}
    Eddy currents are circulating currents induced in the core of an
    inductor by the changing magnetic field. These currents flow in
    closed loops within the core material and can cause significant
    power losses due to the resistive heating of the core. To minimize
    eddy current losses, cores are often laminated or made from
    materials with high electrical resistance, such as ferrites.
  \item \textbf{Types of Inductors}:
    There are many types of inductors. Some common types include:
    \begin{itemize}
      \item \textbf{Air Core Inductors}: No magnetic core, just wire wound
        on a non-magnetic core (or free-standing). No core losses and no
        saturation, but low inductance per unit volume. Used in high-frequency
        and RF applications where core losses would otherwise dominate.
        Example: an air-core coil in a VHF/UHF matching network.
      \item \textbf{Ferrite Core Inductors}: Wound on a ferrite (ceramic
        ferrimagnetic) core. Ferrite's high resistivity keeps eddy current
        losses low even at high frequency, so these are the standard choice
        for switching power supply inductors and RF chokes. \textbf{Ferrites
        saturate at a fairly low flux density}, so core geometry and gap
        must be sized carefully for the expected current.
        Example: the output inductor of a buck converter switching at a few
        hundred kHz.
      \item \textbf{Iron/Steel Core Inductors}: Wound on laminated iron or
        steel. Laminations (thin, insulated sheets) limit eddy current
        losses, but these losses still rise sharply with frequency, so
        \textbf{iron cores are largely restricted to low-frequency, high-power
        applications} (mains frequency and below). They allow much higher
        flux density than ferrites before saturating.
        Example: the core of a 50/60 Hz mains transformer or a large power
        inductor in a linear supply.
      \item \textbf{Toroidal Inductors}: Wound on a doughnut-shaped core
        (can be ferrite, iron powder, or air). The closed magnetic loop
        means the field is largely self-contained, which \textbf{reduces
        electromagnetic interference (EMI) and stray coupling} compared to
        a rod or solenoidal core of similar inductance, and improves
        efficiency. This is really a statement about geometry rather than
        a separate material class.
        Example: a toroidal common-mode choke on a power supply input.
      \item \textbf{Variable Inductors}: Inductance is adjusted mechanically,
        most commonly by moving a ferrite or powdered-iron slug in or out of
        the coil (a slug-tuned inductor), which changes the effective
        permeability seen by the winding. Tapped coils with a switch are
        another common approach.
        Example: a slug-tuned IF coil in an old superheterodyne radio.
    \end{itemize}
    Some special inductors also exist to perform a unique function, such as:
    \begin{itemize}
      \item \textbf{Chokes}: Inductors designed to present high impedance
        to high-frequency signals while passing low-frequency or DC current
        with little attenuation — essentially an inductor used specifically
        for its blocking behaviour rather than for energy storage or
        filtering math.
        Example: a ferrite bead or RF choke in series with a supply line to
        block noise from a switching regulator.
      \item \textbf{Transformers}: Two or more windings magnetically
        coupled through a shared core, used to transfer energy between
        circuits without a direct electrical connection. Important in power
        applications (stepping voltage up or down) and signal applications
        (galvanic isolation, impedance matching).
        Example: a mains transformer stepping 230V down to 12V, or a
        pulse transformer isolating a gate-drive signal.
    \end{itemize}
  \item \textbf{Advise:} These are some important things to keep in mind when
    working with inductors.
    \begin{itemize}
      \item \textbf{Kickback Voltage:} When an inductor is suddenly
        disconnected from a circuit, it \textbf{can generate a sudden and
        short lived (transient) high voltage spike due to the collapsing
        magnetic field}. This phenomenon is known as kickback and
        \textbf{can damage other components} in the circuit. To prevent
        frying your circuit, a diode is often placed across the inductor
        to safely dissipate the energy towards. Other names for this
        phenomenon include flyback voltage, back EMF, and inductive
        kickback. Other names for the diode used to protect are
        flyback, freewheeling, snubber, clamp, catch or suppressor diode.
      \item \textbf{Saturation:} Inductors have a maximum current rating,
        beyond which the \textbf{core material can become saturated and cannot
        increase current proportionally with the applied voltage}. This
        is known as the saturation current. It leads to a loss of inductance
        and energy storage ability. You must ensure that the current through
        the inductor does not exceed this limit.
    \end{itemize}
\end{itemize}

\subsection{Diode}

A diode is a two-terminal, active, non-linear and time-variant
device that allows current to flow in one direction while blocking
it in the opposite direction. The most common type of diode is the
semiconductor diode, which is made from a piece of semiconductor
material (usually silicon) that has been doped with impurities to create
a p-n junction. Other types of diodes include Schottky diodes, Zener
diodes, light-emitting diodes (LEDs) and photodiodes.

\section{Bipolar Junction Transistor (BJT)}

A BJT is three-port, active, non-linear and time-variant device.
It is physically built from two PN junctions that share a single
region, with the two junctions oriented oppositely. This shows
up in its governing model. It is called a \emph{bipolar} device
because both electrons and holes are involved in the conduction
process.

The three terminals are the base (B), the collector (C) and the
emitter (E). It is often called a \emph{current controlled device}
because of its base-collector current relationship. In reality, the
base-emitter voltage controls the collector current.

BJT family has two types of transistors: NPN and PNP.\ They are
(in a sense) mirror images of each other.

\begin{circuitfig}
	% Paths, nodes and wires:
	\node[npn] at (5.5, 7.02){};
	\node[pnp, yscale=-1] at (9, 6.98){};
	\node[shape=rectangle, minimum width=1.465cm, minimum height=0.465cm] at (5.5, 6){} node[anchor=north, align=center, text width=1.077cm, inner sep=6pt] at (5.5, 6.25){Emitter};
	\node[shape=rectangle, minimum width=0.97cm, minimum height=0.465cm] at (4.157, 7.02){} node[anchor=north, align=center, text width=0.582cm, inner sep=6pt] at (4.157, 7.27){Base};
	\node[shape=rectangle, minimum width=0.97cm, minimum height=0.465cm] at (7.657, 6.98){} node[anchor=north, align=center, text width=0.582cm, inner sep=6pt] at (7.657, 7.23){Base};
	\node[shape=rectangle, minimum width=1.465cm, minimum height=0.465cm] at (5.5, 8.04){} node[anchor=north, align=center, text width=1.077cm, inner sep=6pt] at (5.5, 8.29){Collector};
	\node[shape=rectangle, minimum width=1.465cm, minimum height=0.465cm] at (9, 8){} node[anchor=north, align=center, text width=1.077cm, inner sep=6pt] at (9, 8.25){Collector};
	\node[shape=rectangle, minimum width=1.465cm, minimum height=0.465cm] at (9, 5.96){} node[anchor=north, align=center, text width=1.077cm, inner sep=6pt] at (9, 6.21){Emitter};
	\node[shape=rectangle, minimum width=1.965cm, minimum height=0.755cm] at (5.5, 4.855){} node[anchor=north, align=center, text width=1.577cm, inner sep=6pt] at (5.5, 5.25){\Large NPN};
	\node[shape=rectangle, minimum width=1.965cm, minimum height=0.755cm] at (9, 4.855){} node[anchor=north, align=center, text width=1.577cm, inner sep=6pt] at (9, 5.25){\Large PNP};
\end{circuitfig}

The behaviour of a BJT can be approximated well by the below equations.
\begin{itemize}
  \item \textbf{Mathematical Model}:
    There are several mathematical models that describe the behaviour of the
    BJT. Some of the most common ones are:
    \begin{itemize}
      \item \textbf{Ideal Model}: The ideal model assumes that the BJT has infinite
        current gain and zero base current.
        \[
          I_C = \begin{cases}
            0 & \text{if } V_{BE} < V_{th} \\
            \infty & \text{if } V_{BE} > V_{th}
          \end{cases}
        \]
      \item \textbf{Non-linear Model}: The Ebers-Moll model combined with Base-Width Modulation
        Non-Linearity is a more accurate model that describes the behaviour of BJT. Ebers-Moll
        is based on from two coupled Shockley's diode equations.
        \[
          I_C = \left[\underbrace{\alpha_F I_{FS} \left( e^{\frac{V_{BE}}{V_T}} - 1 \right)}_{\text{Forward Contribution}} - \underbrace{\alpha_R I_{RS} \left( e^{\frac{V_{BC}}{V_T}} - 1 \right)}_{\text{Reverse Contribution}}\right] \left( 1 + \frac{V_{CE}}{V_A} \right)
        \]
        \begin{itemize}
          \item \(I_C\) is the collector current
          \item \(I_{FS}\) is the forward saturation current
          \item \(\alpha_F\) is the forward common base current gain
          \item \(V_{BE}\) is the base-emitter voltage
          \item \(V_T\) is the thermal voltage
          \item \(I_{RS}\) is the reverse saturation current
          \item \(\alpha_R\) is the reverse common base current gain
          \item \(V_{BC}\) is the base-collector voltage
          \item \(V_{CE}\) is the collector-emitter voltage
          \item \(V_A\) is the Early voltage
        \end{itemize}
        We \textbf{will make two simplifications to our Non-Linear model} to make
        it easier to work with,
        \begin{itemize}
          \item Ignore the dependence of collector current on collector-emitter
            voltage (\(V_{CE}\)) by assuming an infinite early voltage \(V_A \to \infty\).
            This simplifies the Ebers-Moll model to:
            \[
              I_C = \alpha_F I_{FS} \left( e^{\frac{V_{BE}}{V_T}} - 1 \right) - \alpha_R I_{RS} \left( e^{\frac{V_{BC}}{V_T}} - 1 \right)
            \]
          \item In forward active mode, the base-collector voltage is reverse biased
            (\(V_{BC} \ll 0\)) and thus \( e^{\frac{V_{BC}}{V_T}} \approx 0 \) and the
            Ebers-Moll model simplifies to:
            \begin{align*}
              I_C &= \alpha_F I_{FS} \left( e^{\frac{V_{BE}}{V_T}} - 1 \right) - \alpha_R I_{RS} \left( e^{\frac{V_{BC}}{V_T}} - 1 \right) \\
              I_C &= \alpha_F I_{FS} \left( e^{\frac{V_{BE}}{V_T}} - 1 \right) - \alpha_R I_{RS}(-1) \\
              I_C &= \alpha_F I_{FS} \left( e^{\frac{V_{BE}}{V_T}} - 1 \right) + \alpha_R I_{RS} \\
              I_C &\approx \alpha_F I_{FS} \left( e^{\frac{V_{BE}}{V_T}} - 1 \right)
            \end{align*}
          \item Forward and Reverse Saturation Currents \( I_{FS} \) and \(I_{RS}\)
            are related to each other by the saturation current \(I_{S}\):
            \[
              \alpha_F I_{FS} = \alpha_R I_{RS} = I_S
            \]
            Thus, the collector current model becomes:
            \[
              I_C \approx I_S \left( e^{\frac{V_{BE}}{V_T}} - 1 \right)
            \]
        \end{itemize}
        We \textbf{can also derive the base and emitter currents} from the collector
        current model.
        \begin{itemize}
          \item Base Current (\(I_B\)):
            The base current is given by:
            \[
              I_B = \frac{I_C}{\beta}
            \]
            where \(\beta\) is the current gain of the BJT.
          \item Emitter Current (\(I_E\)):
            The emitter current is given by:
            \begin{align*}
              I_E &= I_C + I_B \\
              I_E &= I_C + \frac{I_C}{\beta} \\
              I_E &= I_C \left( 1 + \frac{1}{\beta} \right)
              I_E \approx I_C \hspace{1cm} \text{[for large \(\beta\)]}
            \end{align*}
        \end{itemize}
      \item \textbf{Linear Model}: The Hybrid-Pi model models the BJT as a
        dependent current source when linearized around a bias point to set
        its transconductance \(g_m\) property.
        \begin{align*}
          I_C = g_m V_{BE} + g_o V_{CE} \\
          I_B = g_\pi V_{BE}
        \end{align*}
        Here, \(g_m\) is the transconductance, \(g_o\) is the
        output conductance and \(g_\pi\) is the input conductance. Below
        schematic shows the simple idea of the hybrid-pi model for BJT.
        \begin{circuitfig}
        	\draw (9.25, 6.75) to[american current source, l_={$g_m V_{BE}$}] (9.25, 5);
        	\draw (6.5, 6.75) to[european resistor, l={$g_\pi$}] (6.5, 5);
        	\draw (10.75, 6.75) to[european resistor, l={$g_o$}] (10.75, 5);
        	\draw (6.5, 5) -- (9.25, 5);
        	\draw (9.25, 5) -- (10.75, 5);
        	\draw (10.75, 6.75) -- (9.25, 6.75);
        	\draw (10.75, 6.75) -- (11.75, 6.75);
        	\draw (6.5, 6.75) -- (5.75, 6.75);
        	\node[shape=rectangle, minimum width=0.965cm, minimum height=0.465cm](N1) at (5.25, 6.75){} node[anchor=center] at (N1.text){$V_B$};
        	\node[shape=rectangle, minimum width=0.965cm, minimum height=0.465cm](N2) at (8.5, 4.5){} node[anchor=center] at (N2.text){$V_E$};
        	\draw (8.5, 5) -| (8.5, 4.75);
        	\node[shape=rectangle, minimum width=0.965cm, minimum height=0.465cm](N3) at (12.25, 6.75){} node[anchor=center] at (N3.text){$V_C$};
        \end{circuitfig}
        \(g_m\) can be calculated as follows:
        \begin{align*}
          g_m &= \frac{dI_C}{dV_{BE}} |_{V_{DS(Q)}} \\
          g_m &= \frac{d}{dV_{BE}} \left[ I_S \left( e^{\frac{V_{BE}}{V_T}} - 1 \right) \right] |_{V_{DS(Q)}} \\
          g_m &= \frac{d}{dV_{BE}} \left[ I_S e^{\frac{V_{BE}}{V_T}} - I_S \right] |_{V_{DS(Q)}} \\
          g_m &= \left[ I_S e^{\frac{V_{BE}}{V_T}} \times \frac{1}{V_T} \right] \\
          g_m &= \frac{I_C}{V_T} \hspace{1cm}
        \end{align*}
        \(g_o\) can be calculated as follows:
        \begin{align*}
          g_o &= \frac{dI_C}{dV_{CE}} |_{V_{BE(Q)}} \\
          g_o &= \frac{d}{dV_{CE}} \left[ I_S \left( e^{\frac{V_{BE}}{V_T}} - 1 \right) \left( 1 + \frac{V_{CE}}{V_A} \right) \right] |_{V_{BE(Q)}} \\
          g_o &= \left[ I_S \left(e^{\frac{V_{BE}}{V_T}} - 1 \right) \times \frac{1}{V_A} \right] \\
          g_o &= \frac{I_C}{V_A} \hspace{1cm}
        \end{align*}
    \end{itemize}
  \item \textbf{Voltages \& Relations}:
    There are several important voltages in a BJT:\
    \begin{itemize}
      \item \textbf{Thermal/Boltzmann Voltage (\(V_T\))}:
        The thermal voltage is given by:
        \[
          V_T = \frac{kT}{q}
        \]
        where \(k\) is the Boltzmann constant, \(T\) is the temperature in
        Kelvin and \(q\) is the charge of an electron. At room temperature,
        \(V_T \approx 26\,\mathrm{mV}\):
      \item \textbf{Threshold Voltage (\(V_{th}\))}:
        The threshold voltage is the minimum voltage required to bias a junction
        such that it conducts. For a silicon junction (most common), the
        threshold voltage is \(V_{th} \approx 0.6\text{--}0.7\,\mathrm{V}\).
        Since we have two junctions in a BJT, we have two threshold voltages.
        We define \(V_{Fth}\) for the forward junction threshold voltage and
        \(V_{Rth}\) for the reverse junction threshold voltage.
    \end{itemize}
    They are also related to each other:
    \begin{itemize}
      \item \textbf{Base-Emitter \(V_{BE}\) and Base-Collector \(V_{BC}\) Voltages}:
        If you rearranged Ebers-Moll model for \(V_{BE}\) or \(V_{BC}\) after
        simplifying for forward active or reverse active mode, you get:
        \begin{itemize}
          \item \textbf{Under Forward Active Mode}:
            \begin{align*}
              V_{BE} = V_T \ln \left( \frac{I_C}{I_{S}} + 1 \right) \\
              V_{BC} \ll 0
            \end{align*}
          \item \textbf{Under Reverse Active Mode}:
            \begin{align*}
              V_{BE} \ll 0 \\
              V_{BC} = V_T \ln \left( \frac{I_C}{I_{S}} + 1 \right)
            \end{align*}
          \end{itemize}
      \item \textbf{Collector-Emitter \(V_{CE}\) Voltage}:
        The collector-emitter voltage is the voltage across the collector and
        emitter terminals. It is given by:
        \begin{align*}
          V_{CE} &= V_{BE} - V_{BC} \\
          V_{CE} &= V_B - V_E - (V_B - V_C) \\
          V_{CE} &= V_B - V_E - V_B + V_C \\
          V_{CE} &= V_C - V_E
        \end{align*}
    \end{itemize}
  \item \textbf{Modes of Operation}:
    A BJT can operate in four different modes depending on \(V_{BE}\) and
    \(V_{CE}\). If \(V_{Fth} \approx 0.7V\), \(V_{Rth} \approx -0.7 V\)
    and \(\beta_R\) is the reverse current gain, then:
    \begin{table}[H]
      \centering
      \begin{tabular}{|c|c|c|c|}\hline
        \textbf{Mode}              & \( V_{BE} \)        & \( V_{BC} \)             & \(I_{C}\) \\ \hline
        Cutoff                     & $V_{BE} < 0.7       $ & $V_{BC} < -0.7 $       & $I_C \approx 0$ \\ \hline
        Forward Active (Amplifier) & $V_{BE} \gtrsim 0.7 $ & $V_{BC} < -0.7 $       & $I_C \approx \beta I_B$ \\ \hline
        Saturation (Switch)        & $V_{BE} \gtrsim 0.7 $ & $V_{BC} \gtrsim -0.7 $ & $I_C \approx I_{C(sat)}$ \\ \hline
        Reverse Active             & $V_{BE} < 0.7       $ & $V_{BC} \gtrsim -0.7 $ & $I_C \approx \beta_R I_B$ \\ \hline
      \end{tabular}
    \end{table}
  \item \textbf{Applications}:
    Various modes can be exploited for different applications. The two
    most common are switching and amplification.
    \begin{itemize}
      \item \textbf{Switching}:
        A BJT can be used as a switch by being driven between \textbf{cutoff}
        (off) and \textbf{saturation} (on). The idea is size the base resistor
        such that there is enough \(I_C\) -- it collapses  \(V_{CE}\) to
        \(V_{CE(sat)}\) (typically \(0.2\text{--}0.3 V\)) through whatever load
        that exists on the collector. You only need to,
        \begin {itemize}
          \item Find the required collector current \(I_C\) to drive the load.
            \[
              I_C = \frac{V_{CC} - V_{CE(sat)}}{R_C}
            \]
          \item Find the required base current \(I_B\) to drive the transistor into
            saturation. \(\beta_F\) is forced current gain to ensure large base current
            and even larger collector current, usually \(\beta_F = 10\).
            \[
              I_B = \frac{I_C}{\beta_{F}}
            \]
          \item Find the base resistor \(R_B\) to limit the base current to \(I_B\).
            \[
              R_B = \frac{V_{IN} - V_{BE}}{I_B}
            \]
        \end{itemize}
      \item \textbf{Amplification}:
        A BJT can be used as an amplifier by being biased in \textbf{forward
        active} mode. The idea is to size the base resistor such that your
        desired \(I_C\) keeps the transistor in forward active mode for the
        entire swing of the input signal. You only need to,
        \begin{itemize}
          \item Find the required \(I_B\) for desired \(I_C\). The forward
            current gain is usually \(\beta = 100\).
            \[
              I_B = \frac{I_C}{\beta}
            \]
          \item Find the base resistor \(R_B\) to limit the base current
            to \(I_B\).
            \[
              R_B = \frac{V_{IN} - V_{BE}}{I_B}
            \]
        \end{itemize}
    \end{itemize}
  \item \textbf{Effects}:
    BJTs are affected by many non-linear effects that are not captured by
    the Ebers-Moll model. Sometimes they cause hard to diagnose problems.
    \begin{itemize}
        \item \textbf{Base Width Modulation}: Increasing \(V_{CE}\) shrinks the base
          width physically. This increases the collector current and is captured
          by the Early voltage (\(V_A\)) and produces the Early Effect.
          \[
            I_C = I_{S} \left( e^{\frac{V_{BE}}{V_T}} - 1 \right) \left( 1 + \frac{V_{CE}}{V_A} \right) \hspace{1cm} \text{[forward active]}
          \]
    \end{itemize}
  \item \textbf{Noise}:
      BJTs are also affected by noise namely, thermal noise, shot noise and flicker
      noise. You can read more about these noises in~\ref{sec:theory_general_noise}.
  \item \textbf{Circuit Analysis}:
    Depending upon Axes of Linearity requirements, You can either use the
    Non-Linear or Linear model to analyse the behaviour of a BJT while
    performing your circuit analysis.
    \begin{itemize}
      \item \textbf{Non-Linear Analysis}:
        Use Ebers-Moll model to analyse the behaviour of a BJT in a circuit.
      \item \textbf{Linear Analysis}:
        Use small-signal model to analyse the behaviour of a BJT in a circuit.
    \end{itemize}
\end{itemize}

\section{Analysis}

Circuit Analysis is the process of determining a circuit's behaviour from its
topology and component values. Analysis is a deterministic process, with a
variety of techniques available to a circuit designer depending on what is
already known and what is being sought.

When analysing a circuit, some characteristics are known and others are the
target of the analysis. For example,
\begin{itemize}
    \item Given a DC input voltage, what is the output voltage of a
      power supply?
    \item Given an AC input signal, what is the frequency response of a
      Low Pass Filter (LPF)?
    \item Given an AC input signal, what is the gain of a transistor
      amplifier?
\end{itemize}
This raises two natural questions:
\begin{enumerate}
    \item How many distinct kinds of analysis technique exist, in
      principle, for any given circuit?
    \item Given a set of known characteristics (AC/DC input, time/frequency
      response), which technique is the appropriate one to apply?
\end{enumerate}
To answer the first question, we build a structure that classifies every
analysis technique. Each characteristic that distinguishes one technique from
another is binary. For example, the input signal is either DC or AC; the
response is computed in either the time domain or the frequency domain; the
non-linear devices in the circuit are either retained in their full
non-linear form or replaced with a first-order (small-signal) approximation
about a bias point. The second question, which technique to apply for a
given circuit, depends on the specific problem at hand; later sections
address this with worked examples.

\subsection{Classification Space}

Figure~\ref{fig:analysis-classes} considers three such characteristics to
build a classification space.
\begin{enumerate}
    \item \textbf{Input Excitation:} The input signal is either DC or AC.
      This forms the Axis of Input Excitation, the X-axis of the
      classification space.
    \item \textbf{Mathematical Representation:} The response is computed
      in either the time domain or the frequency domain. This forms the
      Axis of Mathematical Representation, the Y-axis of the classification
      space. You can flip between the time and frequency domain using a
      Fourier or Laplace Transform.
    \item \textbf{Non-Linear Device Modelling:} The non-linear devices are
      either retained in full (large-signal analysis) or linearized about
      a bias point (small-signal analysis). This forms the Axis of
      Linearity, the Z-axis of the classification space.
\end{enumerate}

\begin{figure}[H]
    \centering
    \includegraphics[width=0.8\textwidth]{01_ANALYSIS_CLASSES.png}
    \caption{Classification space for analysis techniques.}\label{fig:analysis-classes}
\end{figure}

\subsection{Analysis Techniques}

The resulting space has eight quadrants, each corresponding to a distinct
analysis technique. Let us now consider each quadrant in turn, and the
analysis technique that arises from it.

\begin{itemize}
    \item \textbf{Q1: Small Signal, Time Domain, DC Input:}
      The analysis of a circuit with a DC input
      signal, where the response is computed in the time domain, and the
      non-linear devices are linearized about a bias point. In small
      signal analysis, the DC operating point is used to compute the
      linearization and is then removed from the circuit: independent DC
      voltage sources are shorted, independent DC current sources are
      opened, leaving only the incremental (perturbation) behaviour about
      that point. A DC input to this linearized circuit is therefore
      identically zero by construction. This analysis is \textbf{not useful}.
    \item \textbf{Q2: Small Signal, Time Domain, AC Input:}
      The analysis of a circuit with an AC input
      signal, where the response is computed in the time domain, and the
      non-linear devices are linearized about a bias point. It is a very
      useful technique, and is equal to the Inverse Laplace Transform of
      \textbf{Q4}.
    \item \textbf{Q3: Small Signal, Frequency Domain, DC Input:}
      It is the Laplace Transform of \textbf{Q1}, differing only in the Mathematical
      Representation used. Since \textbf{Q1} is identically zero, so is this one, by
      linearity of the transform. This analysis is \textbf{not useful}.
    \item \textbf{Q4: Small Signal, Frequency Domain, AC Input:}
      The analysis of a circuit with an AC
      input signal, where the response is computed in the frequency
      domain, and the non-linear devices are linearized about a bias
      point. It is a very useful technique. In simulators like ngspice, it
      is implemented as AC analysis.
    \item \textbf{Q5: Large Signal, Time Domain, DC Input:}
      The analysis of a circuit with a DC input
      signal, where the response is computed in the time domain, and the
      non-linear devices are retained in full. It is a very useful
      technique. In simulators like ngspice, it is implemented as DC
      analysis, OP analysis, or DC sweep analysis.
    \item \textbf{Q6: Large Signal, Time Domain, AC Input:}
      The analysis of a circuit with an AC input
      signal, where the response is computed in the time domain, and the
      non-linear devices are retained in full. It is a very useful
      technique. In simulators like ngspice, it is implemented as TRAN
      analysis.
    \item \textbf{Q7: Large Signal, Frequency Domain, DC Input:}
      It is the Laplace Transform of \textbf{Q5}, differing only in the
      Mathematical Representation used. It does contain information about
      the circuit, namely the frequency response to a DC input but since
      it is only a single frequency bin at zero, This analysis is
      \textbf{not useful}.
    \item \textbf{Q8: Large Signal, Frequency Domain, AC Input:}
      It is the Laplace Transform of \textbf{Q6},
      differing only in the Mathematical Representation used. It contains
      useful information about the circuit, namely the frequency response
      of the circuit to an AC input with one caveat: unlike
      \textbf{Q4}, the non-linear device models are not
      linearized, so superposition does not hold. The resulting spectrum
      is only valid for the specific input amplitude and waveform used to
      compute it, not a general transfer function; changing the input
      requires recomputing this quadrant from \textbf{Q6},
      rather than scaling the result. In simulators like ngspice, it is
      implemented as FOUR analysis.
\end{itemize}

Removing \textbf{Q1}, \textbf{Q3} and \textbf{Q7} due to lack of usefulness,
We can summarize of the quadrants and the analysis techniques that arise
from them:
\begin{table}[H]
    \centering
    \begin{tabular}{|c|c|c|c|c|c|}\hline
      \textbf{Q} & \textbf{ZYX} & \textbf{Linearity} & \textbf{Representation} & \textbf{Input Excitation} & \textbf{SPICE} \\ \hline
      Q2 & 001 & Small Signal & Time      & AC & None \\ \hline
      Q4 & 011 & Small Signal & Frequency & AC & AC \\ \hline
      Q5 & 100 & Large Signal & Time      & DC & DC, OP, DC Sweep \\ \hline
      Q6 & 101 & Large Signal & Time      & AC & TRAN \\ \hline
      Q8 & 111 & Large Signal & Frequency & AC & FOUR \\ \hline
    \end{tabular}
    \caption{Summary of the eight quadrants of analysis techniques.}\label{tab:quadrant-summary}
\end{table}

\subsection{Quadrant/Axes Dependencies}

Certain Quadrants and Axes are dependent on each other, and this is important to
understand when analysing circuits.
\begin{itemize}
    \item \textbf{Q2} and \textbf{Q4} depend on \textbf{Q5} to find a bias
      point to linearize the non-linear devices about.
    \item \textbf{Q6} can be computed without any knowledge of \textbf{Q5}
      but it is often useful to compute \textbf{Q5} first to find a bias
      point and use that to set initial conditions.
    \item \textbf{An AC Input} can make use of either a small signal or
      large signal device model but it is often convenient to use a small
      signal model, since it is linear and therefore easier to analyse.
      This is done below in the section on hand analysis simplifications.
\end{itemize}

\subsection{Hand Analysis Simplifications}

When analysing circuits by hand, it is often useful to simplify the
circuit schematic by treating various devices as either short circuits
or open circuits, depending on the input excitation AC/DC.\
\begin{table}[H]
    \centering
    \begin{tabular}{|c|c|c|}\hline
      \textbf{Device}   & \textbf{DC Input} & \textbf{AC Input} \\ \hline
      DC Voltage Source & -                 & Short \\ \hline
      AC Voltage Source & Short             & -     \\ \hline
      DC Current Source & -                 & Open  \\ \hline
      AC Current Source & Open              & -     \\ \hline
      Resistor          & -                 & -     \\ \hline
      Capacitor         & Open              & Short (HF/Large C only) \\ \hline
      Inductor          & Short             & Open (HF/Large L only) \\ \hline
      Diode             & -                 & Linearize \\ \hline
      MOSFET            & -                 & Linearize \\ \hline
      BJT               & -                 & Linearize \\ \hline
      Dependent Sources & -                 & -     \\ \hline
      OpAmp             & -                 & -     \\ \hline
    \end{tabular}
    \caption{Hand analysis simplifications for various devices.}\label{tab:hand-analysis-simplifications}
\end{table}

In the above table, \textbf{-} indicates that the device is treated normally
that is without any simplification. \textbf{Short} indicates that the device
should be replaced with a short circuit and \textbf{Open} indicates that the
device should be replaced with an open circuit. \textbf{Linearize} indicates
that the full non-linear model of a device should be replaced with its
small-signal linearized model (about a bias point).

